\section{Security Analysis}

\label{sec:security}

\subsection{Cryptographic Security}

Both TFHE and BFV derive their security from Ring-LWE~\cite{lyubashevsky2010ideal}.
Our parameter selection:

\begin{table}[ht]
\centering
\caption{Security parameters.}
\label{tab:security-params}
\begin{tabular}{lrrr}
\toprule
\textbf{Scheme} & \textbf{$N$} & \textbf{$\log_2 q$} & \textbf{Classical Security} \\
\midrule
TFHE & 1{,}024 & 27 & 128 bits \\
BFV-128 & 4{,}096 & 218 & 128 bits \\
BFV-192 & 8{,}192 & 438 & 192 bits \\
BFV-256 & 16{,}384 & 881 & 256 bits \\
\bottomrule
\end{tabular}
\end{table}

Security levels are estimated using the lattice-estimator~\cite{albrecht2015lattice}
with BKZ-$\beta$ cost model and core-SVP hardness assumption.

\subsection{Semantic Security}

FHE provides IND-CPA (indistinguishable under chosen plaintext attack) security:
an adversary cannot distinguish encryptions of different plaintexts. This means:

\begin{itemize}[leftmargin=2em]
  \item Encrypted balances reveal nothing about the plaintext balance.
  \item Encrypted order prices reveal nothing about the plaintext price.
  \item Encrypted votes reveal nothing about the plaintext vote.
\end{itemize}

FHE does \emph{not} provide IND-CCA2 security (indistinguishable under adaptive
chosen ciphertext attack). Ciphertexts are malleable by design---this is what
enables homomorphic computation. Access control for decryption (Section~\ref{sec:fhe-vm})
is therefore critical: only authorized parties can decrypt, preventing an adversary
from using the decryption oracle to break semantic security.

%% ========================================================================
%%  9. RELATED WORK
%% ========================================================================
